% Source pp. 49-55 (Lectures1-6.pdf).
\section{Path integral for a relativistic particle}

So far, we have discussed the path integral representation in ordinary quantum
mechanics. The real-time transition amplitude is represented as a suitably
defined integral
\[
\bra{q_f}e^{-\frac i\hbar HT}\ket{q_i}=\int\bigl[Dq(t)\bigr]_{\substack{q(0)=q_i\\q(T)=q_f}}\,e^{\frac i\hbar S[q(t)]}\,,
\]
over all ``paths'' $q(t)$ from $q_i$ to $q_f$. Now I want to discuss how the
idea of the path integral can be applied to a relativistic particle.

As we have seen (Lecture~4), the KG theory is the theory of spinless Bose
particles, and we expect that the Feynman propagator $D_F(x_f-x_i)$ (and
related Euclidean correlation function $D_E$) admit an interpretation in terms
of amplitude associated with the particle propagation from $x_i$ to $x_f$. Can
we write such amplitude as an integral over the ``paths''?

In classical mechanics the action of a relativistic particle travelling from
$x_i$ to $x_f$ is
\[
S=-m\int_{x_i}^{x_f}\sqrt{dx_\mu\,dx^\mu}=-m\int_{(t_i,\xx_i)}^{(t_f,\xx_f)}\sqrt{dt^2-d\xx^2}\,,
\]
where the integration is performed along a given path (generally described by
a parametric curve $x^\mu(s)$) from $x_i$ to $x_f$. In classical theory it is
assumed that $(dx^0)^2\ge d\xx^2$ along the whole path (the particle velocity
never exceeds the speed of light). Then the time component $t=x^0$ of $x^\mu$
can serve as the parameter along the path, and we can write
\[
S[\xx(t)]=-m\int_{t_i}^{t_f}dt\,\sqrt{1-\dot\xx^2}\,.
\]
In quantum theory one may attempt to write for the transition amplitude
\[
\int_{\text{paths}(i\to f)}e^{\frac i\hbar S}\,[D\xx(t)]\,,
\]
where the integration is over ``all paths'' from $x_i$ to $x_f$
(figure~\ref{fig:relpath}).
\begin{figure}[ht]
\centering\figfile{relpath}
\caption{Two of the paths from $x_i$ to $x_f$ that enter the sum. The
element $dx^\mu$ of a path may be time-like or space-like.}
\label{fig:relpath}
\end{figure}

The problems appear when we want to give a precise meaning to this expression.
On the formal side, we observe that the Hamiltonian associated with the
problem,
\[
H=\sqrt{m^2+\pp^2}
\]
(where $\pp$ is the momentum conjugated to $\xx$), is not ``local'' in the
sense we discussed the last time; the matrix elements $\bra{\xx}H\ket{\xx'}$
do not vanish when $\xx\neq\xx'$ (check it against the Exercise in
Lecture~2). More fundamentally, the problem is about the class of the paths
which we need to include into the path integral. The expression
$dx_\mu dx^\mu$ in the square root in $S$ can be either positive or negative,
depending on whether $dx^\mu$ is time-like or space-like. One can try to
exclude from the integration the paths having space-like $dx^\mu$ anywhere,
but it does not seem to be right idea if we want to describe the propagator
$D_F$ this way. Indeed, unlike the commutator, the function $D_F(\xx,t)$ does
not vanish outside the light cone (at $|t|<|\xx|$). It is exponentially small
if you go far outside the light cone, but it is not zero. This is similar to a
non-relativistic particle in the classically inaccessible region of its
configuration space: the wave function is exponentially small but does not
vanish (``quantum tunneling''). Thus we have to admit the paths having
space-like $dx^\mu$ somewhere along them. Once such paths are admitted, the
Lorentz invariance requires that paths in which $dx^0$ is allowed to take
negative values (i.e.\ the paths with the folds going ``backward in time'')
are admitted as well (figure~\ref{fig:backpath}).
\begin{figure}[ht]
\centering\figfile{backpath}
\caption{A path with a fold going backward in time.}
\label{fig:backpath}
\end{figure}

For such paths $\xx(t)$ is not a function, and the integral over $[D\xx(t)]$
does not apply. Besides, whereas for the case of space-like $dx^\mu$ there is
a natural choice of the branch of the square root $\sqrt{dx_\mu dx^\mu}$ in the
action ($e^{iS}$ must receive exponentially small contributions from the
space-like parts of the path), we have no clear intuition which sign to choose
in the case of the time-like ``backward in time'' paths.

All these ambiguities are resolved by adopting the Euclidean ($=$ imaginary
time) formulation of the path integral. Replacing as before
\[
t=-i\,x_4\,,
\]
we have $iS\to-\mathcal A$, where the Euclidean action is
\[
\mathcal A=m\int\sqrt{(dx_E)^2}\,,
\]
with explicitly positive expression
\[
dx_E^2=dx_4^2+d\xx^2
\]
under the square root. In the subsequent discussion I will often omit the
subscript $E$ for the Euclidean 4-vectors. With this, one can define the
Euclidean amplitude as the integral
\[
\int_{\text{paths}(x_i\to x_f)}e^{-\mathcal A[\text{path}]}
\]
over all paths connecting two Euclidean points $x_i$ and $x_f$; one expects
that the real-time propagator will be accessible through analytic
continuation, as we previously discussed.

Note that
\[
\mathcal A[\text{path}]=m_0\,L[\text{path}]\,,
\]
where $L[\text{path}]$ is just the Euclidean length of the path. I have
replaced $m$ by the ``bare mass'' $m_0$; we will see that this parameter does
not necessarily coincide with the physical mass but requires
``renormalization''.

By analogy with our definition of the path integral in ordinary quantum
mechanics, let us define this path integral as a limit of finite-dimensional
integrals. Namely, for given (large) number $n$, let us consider the piecewise
linear paths from $x_i=0$ to $x_f=x$, each linear piece having fixed length
$\Delta$ (figure~\ref{fig:steps}).
\begin{figure}[ht]
\centering\figfile{steps}
\caption{A piecewise linear path from $0$ to $x$ made of $n$ pieces of equal
length $\Delta$.}
\label{fig:steps}
\end{figure}
The overall length of such path is
\[
L_n=n\,\Delta\,.
\]
Let $\nu_k$ be the unit 4-vector in the direction of the $k$-th piece;
$|\nu_k|=1$. In this discrete approximation the path integral can be written
as
\[
\sum_{n=0}^\infty e^{-m_0\Delta n}\int\Bigl(\prod_{k=1}^nd\mu(\nu_k)\Bigr)\,\delta^{(4)}\Bigl(\Delta\sum_{k=1}^n\nu_k-x\Bigr)\,.
\]
Here $d\mu(\nu_k)$ stands for the usual measure on $S^3$, and the
delta-function ensures that the vectors $\Delta\,\nu_k$ add up to $x$, i.e.\
the path indeed is from $0$ to $x=x_f-x_i$. Finally, the sum over $n$ must be
performed since we want to include all paths, with arbitrary length $L$.

The delta-function can be represented as
\[
\delta^{(4)}\Bigl(\Delta\sum\nu_k-x\Bigr)=\int\frac{d^4P}{(2\pi)^4}\,e^{iP(x-\Delta\sum\nu_k)}\,.
\]
For fixed $P$ the integration over $\prod d\mu(\nu_k)$ splits into $n$
identical integrals $I(P\Delta)=\int d\mu(\nu)\,e^{-i\Delta\nu P}$,
\[
\prod_{k=1}^n\int d\mu(\nu_k)\,e^{-i\Delta\nu_kP}=\bigl[I(P\Delta)\bigr]^n\,.
\]
Certainly, the integral $I$ can be evaluated explicitly in terms of the
Bessel's function. However, there is no need in its explicit form. We are
going to send $\Delta\to0$ (in which limit a significant contribution comes
from large $n$). Therefore only the behavior of $I$ at small values of its
argument is important. We find
\[
I(P\Delta)=A-A\,\xi\,P^2\Delta^2+O(\Delta^4)\,,
\]
where exact values of the constants $A=\int_{S^3}d\mu(\nu)=2\pi^2$ and
$\xi=1/8$ are not important. We have
\[
I^n\sim A^n\,e^{-\Delta^2n\,\xi P^2}\,.
\]
Now the integration over $P$ can be performed, and we obtain for given $n$
\[
\bigl(4\pi\,n\Delta^2\xi\bigr)^{-2}\,e^{-\frac{x^2}{4\xi n\Delta^2}}\,.
\]
Note that $\Delta$ enters here in the combination $n\Delta^2$, not $n\Delta$
as one could naively expect. Attempting to take the naive limit
\[
\Delta\to0\,,\qquad n\to\infty\qquad\text{with}\quad L=n\Delta=\text{finite}
\]
would lead to the propagation amplitude which vanishes at any finite $|x|$.
This result is not surprising. Our calculation was nearly identical to the
well known probability calculation for a Brownian particle. We know that a
Brownian particle with typical microscopic velocity $\bar v$ will be found,
after the elapsed time $T$, at the distance
\[
\bar x=\bar v\,\sqrt{T\,\Delta t}
\]
away from its initial position. Here $\Delta t$ is the typical time between
collisions. In this analogy
\[
\Delta\sim\bar v\,\Delta t\,,\qquad n\sim T/\Delta t\,.
\]
Thus we expect to have
\[
\bar x^2=n\,\Delta^2\,,
\]
in agreement with our calculation.

We conclude that the right continuous limit is achieved by taking
\[
\Delta\to0\,,\qquad n\to\infty\qquad\text{with}\quad s=\xi\,n\,\Delta^2=\text{finite}\,,
\]
and it is the parameter $s$ that characterizes the geometry of the typical
path; the microscopic length $L=n\Delta$ becomes infinitely large in this
limit. The absolute majority of the paths contributing to the functional
integral are extremely crumpled at all scales, as figure~\ref{fig:crumpled}
illustrates.
\begin{figure}[ht]
\centering\figfile{crumpled}
\caption{A typical path is crumpled at all scales: a small piece of it,
magnified (right), looks as crumpled as the whole.}
\label{fig:crumpled}
\end{figure}
The parameter $s$ has the dimension of the $[\text{length}]^2$; it is said that
the typical path has the fractal dimension 2.

Finally, there are two unpleasant factors to be dealt with before we can take
the limit $\Delta\to0$. These are
\[
A^n\,e^{-m_0n\Delta}=e^{-n(m_0\Delta-\log A)}\,.
\]
Consider $m_0$ as a parameter which can carry some dependence on $\Delta$, and
is free to be adjusted (in the limit $\Delta\to0$) so that
\[
\Delta\,m_0(\Delta)-\log A\ \to\ \xi\,m^2\,\Delta^2\qquad\text{as}\quad\Delta\to0\,,
\]
where it is $m^2$, not $m_0$, which remains constant in this limit. Then the
above two factors combine into
\[
e^{-m^2\xi n\Delta^2}\to e^{-m^2s}\,.
\]
For large $n$ the sum is replaced by the integral over $s$, and one obtains
\[
\int_0^\infty ds\,e^{-m^2s}\int\frac{d^4P}{(2\pi)^4}\,e^{-sP^2}\,e^{iPx}\,,
\]
which is exactly the Schwinger's proper-time representation \eqref{eq:Schwinger}
of the Euclidean correlation function $D_E(x)$ of the KG theory. Thus this
correlation function admits interpretation as the integral over the paths in
Euclidean space-time. As we know, the Feynman propagator $D_F$ coincides with
the analytic continuation of this correlation function back to the real time,
$x_4=it$.

\begin{remark}
In our calculation above we have assumed that the linear pieces of the
discretized paths have equal lengths $\Delta$. This is not essential, and is
done for simplicity only. One can change details of the finite-dimensional
approximations, but in the limit the same correlation function $D_E$ comes
out (though the relation between the ``bare'' and ``renormalized'' mass
parameters $m_0$ and $m$ does depend on the details of the approximation).
This is an elementary example of \emph{universality}.
\end{remark}

In principle, it is possible to develop a theory of interacting relativistic
particles based on the idea of integrating over all the particles paths,
including the paths of many particles, with interaction events (``vertices'').
But much more universal and powerful approach is based on quantization of the
field variables. In the associated path integral, one integrates not over the
paths of the particles, but over the paths of the field variables, i.e.\ over
all the ``trajectories'' $\Phi(\xx,t)$ of the fields. In what follows we
elaborate this idea in the KG theory.
